Nullniveau der Lageenergie: der Punkt, wo der Stein das Gummiband verlässt. Zustand 1: das Band gespannt (nur Spannenergie). Zustand 2: der Stein verlässt das Band (nur kinetische Energie). Zustand 3: im höchsten Punkt (nur Lageenergie).

\begin{abcliste}
\abc
Sei $k$ die Federkonstante, $s$ die Dehnung und $m$ die Masse des Steins (in Kilogramm!). Die Spannenergie im Zustand 1 wird zur kinetischen Energie im Zustand 2:
\begin{align}
 \frac{1}{2}\,k\,s^2 &= \frac{1}{2}\,m\,v^2 \\
 v &= \vF \\
   &= \sG\cdot\sqrt{\frac{\k}{\m}} \\
   &= \v \approx \result{\vP{0}{2}}
\end{align}

\abc
Zuoberst, im Zustand 3, ist alles zu Lageenergie geworden:
\begin{align}
 \frac{1}{2}\,k\,s^2 &= m\,g\,h \\
 h &= \hF \\
   &= \frac{\k\cdot\left(\sG\right)^2}{2\cdot\m\cdot\g} \\
   &= \h \approx \result{\hP{0}{2}}
\end{align}
\end{abcliste}
